\chapter{培养西藏高中生数学运算能力的措施}


\section{注重学生数学思维的培养}

% 培养学生的数学思维就是培养学生思维中的严谨性、简捷性和联想力，能使学生的运算快速且准确。教师在教学过程中由浅入深、层层递进，在设计问题时难度适中，对不同学生的能力进行分层。题目设计一题多解或多题一解，一题多解是锻炼学生的思维能力，多题一解是培养学生对知识的归纳与总结。培养学生良好的数学思维能提升教学工作的效率。

数学思考是在教学中训练学生的严谨、简便性和连贯性，使得他们能够迅速、精确地进行计算。在教学中，老师从浅入深，层层递进，以中等困难的方式来解决问题，并将不同的学生按能力层次划分。一题多解训练了学生的思考，而“多题一题”则是训练了他们对知识的归纳和总结\cite{董安妮关于西藏高中生数学运算能力现状研究}。提高学生的数学思维，提高教学质量。

\begin{example}
  已知 $\overrightarrow{O A}=({k}, 4), \overrightarrow{O B}=(2,5), \overrightarrow{O C}=(3 {k}, 9)$, 且 ${A}, {B}, {C}$ 三点共线, 那么 ${k}= \underline{\hspace{2em}}$.
\end{example}

\begin{proof}[解]
  解法一：距离公式法 \par
  由 $A , B ,C$ 三点共线可得: $|\overrightarrow{AB}|+|\overrightarrow{B C}|=|\overrightarrow{AC}|$, 取 0 点的坐标为 $(0,0)$,
  则 $A(k, 4)$, $B(2,5)$, $C(3 k, 9)$, $\therefore$
  $|\overrightarrow{AB}|=\sqrt{(2-k)^{2}+(5-4)^{2}}$, 
  $|\overrightarrow{B C}|=\sqrt{(3 k-2)^{2}+(9-4)^{2}}$, 
  $|\overrightarrow{AC}|=\sqrt{(3 k-k)^{2}+(9-4)^{2}} $.
  % \begin{align*}
  %   |\overrightarrow{AB}|&=\sqrt{(2-k)^{2}+(5-4)^{2}}, \\
  %   |\overrightarrow{B C}|&=\sqrt{(3 k-2)^{2}+(9-4)^{2}}, \\
  %   |A \vec{C}|&=\sqrt{(3 k-k)^{2}+(9-4)^{2}} 
  % \end{align*}
  $\because|\overrightarrow{AB}|+|\overrightarrow{B C}|=|\overrightarrow{AC}|$ 解得 $k=\frac{10}{7}$.

  解法二：向量共线法 \par
  由 $A, B , C$ 三点共线可得: $\overrightarrow{AB} / / \overrightarrow{B C} / / \overrightarrow{A C}$.
  $\overrightarrow{A B}=\overrightarrow{O B}-\overrightarrow{O A}$,
  $\overrightarrow{B C}=\overrightarrow{O C}-\overrightarrow{O B}$,
  ${\overrightarrow{AB}}=(2-{k}, 1), \overrightarrow{B C}=(4,3 {k}-2)$,
  $\therefore x_{1} y_{2}-x_{2} y_{1}=0$, 即 $(2-k) 4-(3 k-2) 1=8-4 k-3 k+2=10-7 k=0$, 解得 $k=\frac{10}{7}$.
  
  解法三：斜率相等法 \par
  由 A、B、C 三点共线可得:$k_{A B}=k_{B C}=k_{A C}$,
  $\because A({k}, 4), {B}(2,5), {C}(3 {k}, 9)$,
  $\therefore k_{A B}=\frac{5-4}{2-k}=\frac{1}{2-k}, k_{B C}=\frac{9-5}{3 k-2}=\frac{4}{3 k-2}$,
  又 $\because k_{A B}=k_{B C}, \frac{1}{2-k}=\frac{4}{3 k-2}$ 解得 ${k}=\frac{10}{7}$.
\end{proof}



\section{注意学生基础知识的培养}

在中学数学教学过程中，必须重视数学基本公式、公式和定理等方面的教育。在教学中，通过设置问题情境与同学进行交互，使其对数学的探究感、探究感、激活课堂氛围，使新旧知识形成一个完整的知识库。鼓励同学们进行团队协作和探究式的学习，当课堂上遇到不能解答的问题时，请同学们帮助，有困难的同学请老师帮助，老师们归纳总结，指出关键和困难。虽然数学这种事情，光是死记硬背是不行的，但也得给他们留出一些空闲的时候，让他们把这些知识都背下来，然后再去用。但是，掌握数学的要诀并不只是背熟相关的方程式，而是要掌握多种不同的知识，以实现计算目标。运用是以记忆为依据的，而运用则以其方式呈现。在这个前提下，老师们就可以根据问题的顺序来安排问题，比较易被弄糊涂的部分，来检验学生的学业表现。


比如在基本不等式 $a+b \geq 2 \sqrt{a b}$ 中, 如果 $a=b$, 则表示等号是有效的。当你使用 基本不等式来求解最值的问题时, 要牢记一正、二定、三相等, 和定积最大, 积 定和最小\cite{付麦霞2013利用基本不等式求最值问题探究} 。(1) 如果两个正数值之和是固定的, 则其乘积是最大的, 也就是 说, 如果 $a>0, b>0$, 则最大为 $M$, 如果 $a=b$, 则等式成立; (2) 如果两个正 数值乘以一定的数值, 则其和具有最小值, 也就是说, 如果 $a>0, b>0$, 则有一 个最小数值 $\mathrm{N}$, 如果 $a=b$, 则等式成立。在学习过程中, 很容易忽视数学原理 的应用, 从而导致计算中出错。要多做几个相同的问题, 使同学们能经常留意到 在他们的对等的建立情况。

\section{注重日常练习规范性的培养}

我们必须不断练习数学，弥补日常练习中的不足，这样我们才能最终赢得高考。首先，老师要做一个模范，做一个模范，孔子说：“其身正，不令而行；其身不正，虽令不从。”许多同学在解答问题时都是模仿老师的思维和练习。所以，在教学中，老师要注意写作的正确性和正确性。其次，在日常复习中，同学们要认真复习，很多题目都是因为没有做好，所以同学们可以在试卷上用笔尖勾勒一些关键的东西。此外，要小心草稿的运用。在进行计算时，需要用到草稿纸，在使用时要做好计划，标明编号，并严格遵循正确的操作程序，尽量做到与考卷一模一样，以免因为抄错而丢了分数。一次的考试，一次就是数千人的差距，千万别为了这一点而错失了自己想要的学校。在完成作业时，要培养学生进行核对的良好习惯。首先检查题目，看看有没有疏漏或错误；然后，他又看了一遍草稿，看看计算有没有错，有没有抄错。平时要经常提醒和督促学生，使他们形成一个好的学习方式，这样不但能提高学生的数学等学科水平，而且能提高他们的自信。




\section{注重教师课堂创新培养}

教师应注重“以人为本”，而非“一言堂”，充分发挥学生的积极性。
在教学过程中，教师要给学生“留白”，让他们提问，让他们当老师来解答他们的疑问。通过在黑板上的讲解，可以让同学的思维和运算能力得到全面的展现，老师可以在老师的指导下，纠正他们的失误，防止二次犯错，并且提高写作水平。教师要重视教学实践，能够使教师对本课程的理解和理解，从而为指导教师接下来的教学做好准备。在教学中，要注意运用现代科技手段，如：积极地研究、创新、运用新的科技手段，增强课堂的趣味。

比如，在三角方程的教学过程中，有很多的三角方程，而在记住的引子里，有一段咒语：“奇形不变，象征形”。老师将口诀解释为：一全正、二正弦、三正切、四余弦\cite{袁蕾2018信息化背景下小组合作式教学的教学设计}；也就是说，在第一个象限里，这个三角形是一个正负的。但是，这主要是在语言层面进行的，而且，没有真正的了解，仍然会发生计算失误。因此，在教学中，老师可以通过运用多种形式展示三角函数图形的变化，从而使其更好地认识到它的实质。



\section{注重良好教学环境的培养}

中学还可以说是又一个新的数学开始。当西藏中学的第一节数学课开始，老师们就会告诉他们，他们的数学有多困难，会让他们有一种抗拒的感觉，他们只要把他们的注意力集中到他们身上就行了。高一的课程虽然很容易，但却很关键，在高中的时候占据了很大的比例，所以高中一年级的数学是一个很好的开始，这对于高一的学生来说，是一个非常关键的阶段。这时，老师可以把日常的事例和数学结合在一起，使他们了解到，数学并非是抽象化的，而是与我们的日常活动紧密相连，增强了对数学的兴趣。到了二年级，学生们开始有了自己的数学学习方式，现在的他们，虽然是在艰难的状态下，但是他们的计算和学习都变得更加的有压力，这是为了迎接即将到来的高中阶段的考试。高三是为高考做准备的高三，在这个时候，他们的算术和计算技能都比较成熟，以便能够考上自己喜欢的学校。老师要在这个时候，为同学们提供一些技巧方面的知识，让他们建立自信，减少学习的负担，让他们更容易应付考试。